\section{Experimental setup}
\textbf{Hardware.} The simulation experiments were conducted on a single local machine equipped with a multi-core CPU. Because the execution environment was a single workstation rather than a distributed HPC cluster, all MPI "processes" were mapped to individual CPU cores on the same node. Consequently, the 2D Cartesian topology described in Section 4 is purely logical, and the inter-process communication occurs entirely via shared memory rather than a physical network interconnect.

\textbf{Inputs.} The initial state of the forest was generated using a fixed random seed to ensure reproducibility across different runs. The grid was populated with a 60\% probability for empty cells and a 40\% probability for trees. The biological constants were fixed at $p = 0.005$ for tree growth and $l = 0.0001$ for spontaneous lightning. We tested the scalability by simulating global grid sizes ranging from $N = 500$ to $N = 2000$.

\textbf{Methodology.} For each configuration (varying the global grid size $N$ and the number of MPI processes $p$), the execution time was measured using the \texttt{time.time()} Python function. To guarantee the accuracy of our temporal metrics, we placed an \texttt{MPI\_Barrier} immediately before starting the timer and another one immediately after the main simulation loop. This strict synchronization ensures that the recorded time encompasses both the computational workload and the collective overhead of the halo exchanges without race condition artifacts.